% GENERATED FILE. Edit canonical lessons, then run npm run generate:books.
% canonical-source-hash: fnv1a64:6edc0f2327794f98
% canonical-lessons: UR-C509-chhachh, UR-C509-bilona, UR-C509-jamana, UR-C509-handi, UR-C509-chhalkana, UR-R509-review-502-509

\chapter{Milk and Curd}
\label{ch:ur-milk-and-curd}
\hlchaptermodality{509}

\begin{chapteropening}
\textbf{By the end of this chapter:} \emph{I can talk about churning, setting curd and spills.}

Complete the last lesson of chapter 509: I can talk about churning, setting curd and spills.
\end{chapteropening}

\section[\texorpdfstring{\ur{چھاچھ} (chhāchh)}{چھاچھ (chhāchh)}]{\ur{چھاچھ} (chhāchh) — buttermilk}

\label{lesson:UR-C509-chhachh}



\begin{quote}
Before the new one: say the Urdu for to cover, then the Urdu for a skimmer.
\end{quote}

\subsection*{You'll want to know: \ur{چھاچھ}}
\textbf{\ur{چھاچھ}} — \emph{chhāchh} — \textquotedblleft{}buttermilk\textquotedblright{}.

\begin{quote}
\ur{گرمیوں} \ur{میں} \ur{چھاچھ} \ur{اچھی} \ur{لگتی} \ur{ہے۔}
\end{quote}

\begin{quote}
\emph{garmiyoṅ meṅ chhāchh achchhī lagtī hai.}
\end{quote}

\begin{quote}
Buttermilk is lovely in the summer.
\end{quote}

Feminine.

\begin{cousinweb}
\textbf{\ur{چھاچھ}} (\emph{chhāchh}) is inherited from \emph{*chācchī}, \textquotedblleft{}buttermilk\textquotedblright{}, a form CDIAL reconstructs from its descendants (CDIAL 5012).
\end{cousinweb}

\begin{grammarlens}[title={What you've built}]
One more word for describing what you lived and felt.
\end{grammarlens}

\subsection*{Your turn}
\begin{itemize}
\raggedright
  \item \emph{Say it:} \emph{chhāchh}
  \item \emph{Say it:} \emph{chhāchh}, once more
  \item \emph{Say it:} \emph{garmiyoṅ meṅ chhāchh achchhī lagtī hai.} — Buttermilk is lovely in the summer.
  \item \emph{Recall:} say the Urdu for to cover, then the Urdu for a skimmer, then say \emph{chhāchh} again
\end{itemize}

\begin{tcolorbox}[breakable,colback=teal!4,colframe=teal!35!black,title={Before you move on}]
What does \ur{چھاچھ} mean? (\textquotedblleft{}Buttermilk\textquotedblright{}.) Say it once more.
\end{tcolorbox}
\section[\texorpdfstring{\ur{بلونا} (bilonā)}{بلونا (bilonā)}]{\ur{بلونا} (bilonā) — to churn}

\label{lesson:UR-C509-bilona}



\begin{quote}
Before the new one: say the Urdu for a skimmer, then the Urdu for buttermilk.
\end{quote}

\subsection*{You'll want to know: \ur{بلونا}}
\textbf{\ur{بلونا}} — \emph{bilonā} — \textquotedblleft{}to churn\textquotedblright{}.

\begin{quote}
\ur{دادی} \ur{نے} \ur{مٹکے} \ur{میں} \ur{دہی} \ur{بلویا۔}
\end{quote}

\begin{quote}
\emph{dādī ne maṭke meṅ dahī biloyā.}
\end{quote}

\begin{quote}
Grandmother churned the yoghurt in a clay pot.
\end{quote}

Past, with \ur{نے} (he or she …): \ur{اس} \ur{نے} \ur{دودھ} \ur{بلویا}, \emph{us ne dūdh biloyā}.

\begin{cousinweb}
\textbf{\ur{بلونا}} (\emph{bilonā}) is inherited: CDIAL lists the Hindi cousin \emph{bilonā}, \textquotedblleft{}to churn\textquotedblright{}, under \emph{*virapati}, a form it reconstructs (CDIAL 11844).
\end{cousinweb}

\begin{grammarlens}[title={What you've built}]
One more word for describing what you lived and felt.
\end{grammarlens}

\subsection*{Your turn}
\begin{itemize}
\raggedright
  \item \emph{Say it:} \emph{bilonā}
  \item \emph{Say it:} \emph{bilonā}, once more
  \item \emph{Say it:} \emph{dādī ne maṭke meṅ dahī biloyā.} — Grandmother churned the yoghurt in a clay pot.
  \item \emph{Recall:} say the Urdu for a skimmer, then the Urdu for buttermilk, then say \emph{bilonā} again
\end{itemize}

\begin{tcolorbox}[breakable,colback=teal!4,colframe=teal!35!black,title={Before you move on}]
What does \ur{بلونا} mean? (\textquotedblleft{}To churn\textquotedblright{}.) Say it once more.
\end{tcolorbox}
\section[\texorpdfstring{\ur{جمانا} (jamānā)}{جمانا (jamānā)}]{\ur{جمانا} (jamānā) — to set (curd)}

\label{lesson:UR-C509-jamana}



\begin{quote}
Before the new one: say the Urdu for buttermilk, then the Urdu for to churn.
\end{quote}

\subsection*{You'll want to know: \ur{جمانا}}
\textbf{\ur{جمانا}} — \emph{jamānā} — \textquotedblleft{}to set (curd)\textquotedblright{}.

\begin{quote}
\ur{امی} \ur{رات} \ur{کو} \ur{دہی} \ur{جماتی} \ur{ہیں۔}
\end{quote}

\begin{quote}
\emph{ammī rāt ko dahī jamātī haiṅ.}
\end{quote}

\begin{quote}
Mum sets the yoghurt at night.
\end{quote}

Past, with \ur{نے} (he or she …): \ur{اس} \ur{نے} \ur{دہی} \ur{جمایا}, \emph{us ne dahī jamāyā}.

\begin{cousinweb}
\textbf{\ur{جمانا}} (\emph{jamānā}) is inherited: CDIAL lists Hindi \emph{jamānā} and \emph{jāman}, \textquotedblleft{}a piece of curd for curdling milk\textquotedblright{}, under Sanskrit \emph{yamyatē}, \textquotedblleft{}is held\textquotedblright{} (CDIAL 10428).
\end{cousinweb}

\begin{grammarlens}[title={What you've built}]
One more word for describing what you lived and felt.
\end{grammarlens}

\subsection*{Your turn}
\begin{itemize}
\raggedright
  \item \emph{Say it:} \emph{jamānā}
  \item \emph{Say it:} \emph{jamānā}, once more
  \item \emph{Say it:} \emph{ammī rāt ko dahī jamātī haiṅ.} — Mum sets the yoghurt at night.
  \item \emph{Recall:} say the Urdu for buttermilk, then the Urdu for to churn, then say \emph{jamānā} again
\end{itemize}

\begin{tcolorbox}[breakable,colback=teal!4,colframe=teal!35!black,title={Before you move on}]
What does \ur{جمانا} mean? (\textquotedblleft{}To set (curd)\textquotedblright{}.) Say it once more.
\end{tcolorbox}
\section[\texorpdfstring{\ur{ہانڈی} (hāṅḍī)}{ہانڈی (hāṅḍī)}]{\ur{ہانڈی} (hāṅḍī) — an earthen cooking pot}

\label{lesson:UR-C509-handi}



\begin{quote}
Before the new one: say the Urdu for to churn, then the Urdu for to set (curd).
\end{quote}

\subsection*{You'll want to know: \ur{ہانڈی}}
\textbf{\ur{ہانڈی}} — \emph{hāṅḍī} — \textquotedblleft{}an earthen cooking pot\textquotedblright{}.

\begin{quote}
\ur{دادی} \ur{مٹی} \ur{کی} \ur{ہانڈی} \ur{میں} \ur{دال} \ur{پکاتی} \ur{تھیں۔}
\end{quote}

\begin{quote}
\emph{dādī miṭṭī kī hāṅḍī meṅ dāl pakātī thīṅ.}
\end{quote}

\begin{quote}
Grandmother used to cook lentils in a clay pot.
\end{quote}

Feminine.

\begin{cousinweb}
\textbf{\ur{ہانڈی}} (\emph{hāṅḍī}) is inherited: CDIAL lists the Hindi cousin \emph{hā̃ḍī}, \textquotedblleft{}an earthen cooking pot\textquotedblright{}, under \emph{*hāṇḍa}, a form it reconstructs (CDIAL 14050).
\end{cousinweb}

\begin{grammarlens}[title={What you've built}]
One more word for describing what you lived and felt.
\end{grammarlens}

\subsection*{Your turn}
\begin{itemize}
\raggedright
  \item \emph{Say it:} \emph{hāṅḍī}
  \item \emph{Say it:} \emph{hāṅḍī}, once more
  \item \emph{Say it:} \emph{dādī miṭṭī kī hāṅḍī meṅ dāl pakātī thīṅ.} — Grandmother used to cook lentils in a clay pot.
  \item \emph{Recall:} say the Urdu for to churn, then the Urdu for to set (curd), then say \emph{hāṅḍī} again
\end{itemize}

\begin{tcolorbox}[breakable,colback=teal!4,colframe=teal!35!black,title={Before you move on}]
What does \ur{ہانڈی} mean? (\textquotedblleft{}An earthen cooking pot\textquotedblright{}.) Say it once more.
\end{tcolorbox}
\section[\texorpdfstring{\ur{چھلکانا} (chhalkānā)}{چھلکانا (chhalkānā)}]{\ur{چھلکانا} (chhalkānā) — to spill (a liquid)}

\label{lesson:UR-C509-chhalkana}



\begin{quote}
Before the new one: say the Urdu for to set (curd), then the Urdu for an earthen cooking pot.
\end{quote}

\subsection*{You'll want to know: \ur{چھلکانا}}
\textbf{\ur{چھلکانا}} — \emph{chhalkānā} — \textquotedblleft{}to spill (a liquid)\textquotedblright{}.

\begin{quote}
\ur{بچے} \ur{نے} \ur{میز} \ur{پر} \ur{دودھ} \ur{چھلکایا۔}
\end{quote}

\begin{quote}
\emph{bachche ne mez par dūdh chhalkāyā.}
\end{quote}

\begin{quote}
The child spilt milk on the table.
\end{quote}

Past, with \ur{نے} (he or she …): \ur{اس} \ur{نے} \ur{پانی} \ur{چھلکایا}, \emph{us ne pānī chhalkāyā}.

\begin{cousinweb}
\textbf{\ur{چھلکانا}} (\emph{chhalkānā}) is the causative of \emph{chhalaknā}, \textquotedblleft{}to be spilt\textquotedblright{}, which CDIAL lists under \emph{*chala}, a form it reconstructs (CDIAL 5002).
\end{cousinweb}

\begin{grammarlens}[title={What you've built}]
One more word for describing what you lived and felt.
\end{grammarlens}

\subsection*{Your turn}
\begin{itemize}
\raggedright
  \item \emph{Say it:} \emph{chhalkānā}
  \item \emph{Say it:} \emph{chhalkānā}, once more
  \item \emph{Say it:} \emph{bachche ne mez par dūdh chhalkāyā.} — The child spilt milk on the table.
  \item \emph{Recall:} say the Urdu for to set (curd), then the Urdu for an earthen cooking pot, then say \emph{chhalkānā} again
\end{itemize}

\begin{tcolorbox}[breakable,colback=teal!4,colframe=teal!35!black,title={Before you move on}]
What does \ur{چھلکانا} mean? (\textquotedblleft{}To spill (a liquid)\textquotedblright{}.) Say it once more.
\end{tcolorbox}
\section[(dialogue)]{Looking back: from Rattles and Rust to Milk and Curd}

\label{lesson:UR-R509-review-502-509}



\begin{quote}
Say the Urdu for to rattle and a hinge.
\end{quote}

\subsection*{Your turn}
Say these aloud:

\begin{itemize}
\raggedright
  \item to rattle, a hinge, a nut (for a bolt), rust, a metal spring
  \item damp (in a wall), a brick, plaster (on a wall), to have (something) painted, to smear (with mud)
  \item a switch, a plug (on a cable), electric current, a peg (on a wall), a socket (for a plug)
  \item waste paper, a cobweb, a loft, to tear (something), to have (something) put away
  \item cleaning, to scour (pots), to stick (something) on, to be washed, to rinse out
  \item a board for rolling dough, to knead, to roll out (dough), to pat (into shape), the heat of a flame
  \item to temper (with hot spiced oil), to roast in a pan, a lid, to cover, a skimmer
  \item buttermilk, to churn, to set (curd), an earthen cooking pot, to spill
\end{itemize}

\begin{tcolorbox}[breakable,colback=teal!4,colframe=teal!35!black,title={Before you move on}]
To rattle? (\textbf{\ur{کھڑکھڑانا}}, \emph{khaṛkhaṛānā}.) A hinge? (\textbf{\ur{قبضہ}}, \emph{qabzā}.) A nut (for a bolt)? (\textbf{\ur{نٹ}}, \emph{naṭ}.) Rust? (\textbf{\ur{زنگ}}, \emph{zang}.) A metal spring? (\textbf{\ur{اسپرنگ}}, \emph{ispring}.) Damp (in a wall)? (\textbf{\ur{سیلن}}, \emph{sīlan}.) A brick? (\textbf{\ur{اینٹ}}, \emph{īṅṭ}.) Plaster (on a wall)? (\textbf{\ur{پلستر}}, \emph{palastar}.) To have (something) painted? (\textbf{\ur{رنگوانا}}, \emph{rangvānā}.) To smear (with mud)? (\textbf{\ur{لیپنا}}, \emph{līpnā}.) A switch? (\textbf{\ur{سوئچ}}, \emph{svich}.) A plug (on a cable)? (\textbf{\ur{پلگ}}, \emph{plag}.) Electric current? (\textbf{\ur{کرنٹ}}, \emph{karanṭ}.) A peg (on a wall)? (\textbf{\ur{کھونٹی}}, \emph{khūṅṭī}.) A socket (for a plug)? (\textbf{\ur{ساکٹ}}, \emph{sākeṭ}.) Waste paper? (\textbf{\ur{ردی}}, \emph{raddī}.) A cobweb? (\textbf{\ur{جالا}}, \emph{jālā}.) A loft? (\textbf{\ur{اٹاری}}, \emph{aṭārī}.) To tear (something)? (\textbf{\ur{پھاڑنا}}, \emph{phāṛnā}.) To have (something) put away? (\textbf{\ur{رکھوانا}}, \emph{rakhvānā}.) Cleaning? (\textbf{\ur{صفائی}}, \emph{safāī}.) To scour (pots)? (\textbf{\ur{مانجھنا}}, \emph{māṅjhnā}.) To stick (something) on? (\textbf{\ur{چپکانا}}, \emph{chipkānā}.) To be washed? (\textbf{\ur{دھلنا}}, \emph{dhulnā}.) To rinse out? (\textbf{\ur{کھنگالنا}}, \emph{khaṅgālnā}.) A board for rolling dough? (\textbf{\ur{چکلا}}, \emph{chaklā}.) To knead? (\textbf{\ur{گوندھنا}}, \emph{gūṅdhnā}.) To roll out (dough)? (\textbf{\ur{بیلنا}}, \emph{belnā}.) To pat (into shape)? (\textbf{\ur{تھاپنا}}, \emph{thāpnā}.) The heat of a flame? (\textbf{\ur{آنچ}}, \emph{āṅch}.) To temper (with hot spiced oil)? (\textbf{\ur{بگھارنا}}, \emph{baghārnā}.) To roast in a pan? (\textbf{\ur{بھوننا}}, \emph{bhūnnā}.) A lid? (\textbf{\ur{ڈھکن}}, \emph{ḍhakkan}.) To cover? (\textbf{\ur{ڈھانپنا}}, \emph{ḍhāṅpnā}.) A skimmer? (\textbf{\ur{کفگیر}}, \emph{kafgīr}.) Buttermilk? (\textbf{\ur{چھاچھ}}, \emph{chhāchh}.) To churn? (\textbf{\ur{بلونا}}, \emph{bilonā}.) To set (curd)? (\textbf{\ur{جمانا}}, \emph{jamānā}.) An earthen cooking pot? (\textbf{\ur{ہانڈی}}, \emph{hāṅḍī}.) To spill? (\textbf{\ur{چھلکانا}}, \emph{chhalkānā}.)
\end{tcolorbox}
